\documentclass[11pt,letterpaper]{article}

\usepackage[margin=1in]{geometry}
\usepackage{amsmath,amssymb}
\usepackage{graphicx}
\usepackage{booktabs}
\usepackage{hyperref}
\usepackage{listings}
\usepackage{xcolor}
\usepackage{microtype}
\usepackage{parskip}
\usepackage{caption}
\usepackage{subcaption}
\usepackage{enumitem}

\hypersetup{
  colorlinks=true,
  linkcolor=blue!60!black,
  urlcolor=blue!60!black,
  citecolor=blue!60!black,
}

\lstset{
  language=Python,
  basicstyle=\ttfamily\small,
  keywordstyle=\color{blue!70!black}\bfseries,
  commentstyle=\color{gray},
  stringstyle=\color{green!50!black},
  showstringspaces=false,
  breaklines=true,
  frame=single,
  rulecolor=\color{black!20},
  backgroundcolor=\color{black!3},
  xleftmargin=1em,
  xrightmargin=1em,
}

\title{\textbf{infeNMPC: An Agentic Engineering Case Study in\\
Infinite-Horizon Economic Nonlinear MPC}}
\author{Owen Gerdes\\
Carnegie Mellon University}
\date{Spring 2026}

\begin{document}
\maketitle

\begin{abstract}
\noindent
\textbf{infeNMPC} is a Python framework for closed-loop Nonlinear Model Predictive Control (NMPC) that approximates the infinite-horizon cost-to-go via a time-compressed terminal integral. It is built on Pyomo and IDAES, supports both economic and tracking objectives, and provides an optional Lyapunov stability constraint. The project was developed entirely through agentic engineering with Claude, evolving from an early prototype to a polished, tested, and documented package over one semester. This report describes the problem the software solves, the key implementation decisions, representative results, and a reflection on how agentic engineering shaped the development process.
\end{abstract}

\tableofcontents
\newpage

%% ============================================================
\section{Introduction}
%% ============================================================

Model Predictive Control (MPC) is the dominant advanced process control strategy in the chemical process industry. At each sample time, MPC solves an optimization problem over a finite prediction horizon, applies the first optimal input, then repeats. The finite horizon is a practical necessity---solving over an infinite future is intractable---but it creates two well-known deficiencies:

\begin{enumerate}[leftmargin=2em]
  \item \textbf{Stability is not guaranteed} unless a terminal constraint or terminal cost function is added. Without this, the controller can be myopic and drive the plant to a limit cycle or infeasibility.
  \item \textbf{Economic optimality is sacrificed} when only a quadratic tracking cost is used. Real chemical processes have nonlinear economics (throughput, energy, product purity tradeoffs) that quadratic costs cannot capture.
\end{enumerate}

Infinite-horizon MPC addresses both issues by including the full cost-to-go in the objective. The challenge is that $\int_0^\infty L(x,u)\,dt$ is infinite-dimensional. The key insight exploited by this project is that the \emph{tanh time transformation}
\[
  \tau = \tanh(\gamma t) \in [0,1)
\]
compresses the physical time axis $[0,\infty)$ onto the bounded domain $[0,1)$. The transformed ODE can then be discretized with a finite number of collocation points, giving a finite-dimensional approximation of the infinite integral. The parameter $\gamma$ is automatically selected so that the first collocation point $\tau_1$ aligns with the transformation: $\gamma = \operatorname{atanh}(\tau_1)/T_s$.

\textbf{infeNMPC} implements this idea as a reusable Python framework. A user provides a Pyomo model describing their process dynamics and objective function; the framework handles all discretization, block linking, IC solving, closed-loop shifting, and data I/O. An optional Lyapunov stability constraint
\[
  V_k \leq V_{k-1} - \delta \cdot L_{k-1}
\]
can be enforced, providing a formal decrease certificate at each MPC iteration.

%% ============================================================
\section{System Architecture}
%% ============================================================

The framework is organized as an installable Python package with the following top-level structure:

\begin{center}
\begin{tabular}{ll}
\toprule
\textbf{Module} & \textbf{Role} \\
\midrule
\texttt{infNMPC\_options.py} & \texttt{Options} dataclass --- all configuration lives here \\
\texttt{make\_model.py} & Pyomo model construction (steady-state, finite, infinite blocks) \\
\texttt{controllers.py} & \texttt{InfiniteHorizonController}, \texttt{FiniteHorizonController} \\
\texttt{plant.py} & \texttt{Plant} --- single-step dynamic simulator \\
\texttt{run\_MPC.py} & \texttt{mpc\_loop()} --- closed-loop orchestrator \\
\texttt{check\_model.py} & Three-level model validator with CLI \\
\texttt{data\_save\_and\_plot.py} & CSV export, figure generation, standalone plot script \\
\texttt{tools/} & Collocation, indexing, initialization, solver, and debug utilities \\
\midrule
\texttt{examples/} & Worked examples: CSTR, binary distillation, ternary distillation, pendulum \\
\texttt{tests/} & 121-test pytest suite \\
\bottomrule
\end{tabular}
\end{center}

\subsection{Closed-Loop Simulation Flow}

Each MPC iteration follows the standard receding-horizon pattern:

\begin{enumerate}[leftmargin=2em]
  \item \textbf{Controller solve:} IPOPT minimizes the combined finite-horizon stage cost plus the infinite-horizon terminal cost integral over the transformed domain.
  \item \textbf{Lyapunov update:} If \texttt{lyap\_flag=True}, the current $V_k = \phi_\text{track}^{\text{finite}} + \beta \cdot \phi_\text{track}^{\text{infinite}}[\tau=1]$ is computed and stored as $V_\text{prev}$ for the next iteration's constraint.
  \item \textbf{Plant warm-start:} The controller's optimal first-FE trajectory is loaded into the plant. MVs are temporarily unfixed, data is loaded, and they are refixed---necessary because Pyomo's \texttt{load\_data} silently skips fixed variables.
  \item \textbf{Plant solve:} IPOPT integrates the plant ODEs forward one sampling interval.
  \item \textbf{IC shift:} The plant's endpoint state is loaded into the controller as the new initial condition via an unfix$\to$load$\to$refix sequence. The controller horizon is shifted by \texttt{shift\_values\_by\_time}.
\end{enumerate}

\subsection{NLP Block Structure}

The infinite-horizon controller model is a \texttt{ConcreteModel} containing two sub-blocks:

\begin{itemize}[leftmargin=2em]
  \item \textbf{\texttt{finite\_block}}: LAGRANGE-RADAU collocation over $[0, N_\text{fe} \cdot T_s]$. Contains the user's state variables, MVs, ODEs, and the stage cost sum.
  \item \textbf{\texttt{infinite\_block}}: LAGRANGE-LEGENDRE collocation over $\tau \in [0,1]$, representing the physical horizon $[T_s, \infty)$. The derivative substitution $\frac{dx}{d\tau} = \frac{T_s}{\gamma(1-\tau^2)}\frac{dx}{dt}$ converts the physical ODE to the transformed domain. Contains the terminal cost ODE $\frac{d\phi}{d\tau} = \frac{\gamma}{T_s}(1-\tau^2)\cdot[L(x,u) - L^\star]$.
\end{itemize}

The two blocks are linked by state equality constraints at their shared interface: $x_{\text{finite}}[t_\text{end}] = x_{\text{infinite}}[\tau=0]$.

%% ============================================================
\section{Key Implementation Decisions}
%% ============================================================

\subsection{Pyomo and IDAES as the Modeling Backend}

The choice of Pyomo over alternatives (CasADi, JuMP, GAMS) was driven by the IDAES process systems engineering ecosystem, which provides:
\begin{itemize}[leftmargin=2em]
  \item A DAE collocation transformation (\texttt{dae.collocation}) compatible with the block structure we needed.
  \item The \texttt{DynamicModelInterface} from \texttt{pyomo.contrib.mpc}, which handles warm-start data loading and time-shifting across indexed Pyomo variables.
  \item Pre-built IPOPT binaries with MA57 (a high-performance sparse direct linear solver) accessible via \texttt{idaes get-extensions}.
\end{itemize}

Pyomo's algebraic-modeling approach also made it possible to write user model files that are readable Python---the model interface is just two functions that add components to a \texttt{ConcreteModel}---rather than symbolic expressions in a domain-specific language.

\subsection{Forked Pyomo for \texttt{clean\_model='delete'}}

Standard Pyomo's DAE collocation transform retains placeholder \texttt{VarData} and \texttt{ConstraintData} entries at non-collocation points (e.g., $t=0$ for LAGRANGE-RADAU). These stale entries caused incorrect IC propagation and warm-start failures: when \texttt{shift\_values\_by\_time} moved data forward, the $t=0$ entries (no longer connected to physics) remained at their previous values and corrupted the next iteration's initial condition.

The fix required adding a \texttt{clean\_model='delete'} option to Pyomo's collocation transform, which removes all non-collocation VarData and ConstraintData after discretization. This change was implemented in a fork (\href{https://github.com/oag00002/pyomo}{oag00002/pyomo}, branch \texttt{mpc-rewrite}) and has been the foundation of all framework stability.

\subsection{The t=0 Constraint Deletion Invariant}

Even with the forked Pyomo, \texttt{equations\_write} must be called \emph{after} discretization (Pyomo's DAE dictionary expansion only works once \texttt{m.time} is fully populated). This means every \texttt{Constraint(m.time, rule=...)} call in \texttt{equations\_write} creates an entry at $t=0$---even though $t=0$ is not a RADAU collocation point and has no physics there.

For the infinite block (LAGRANGE-LEGENDRE), both $\tau=0$ and $\tau=1$ are non-collocation endpoints that receive spurious constraint entries.

The framework deletes these entries immediately after calling \texttt{equations\_write}:

\begin{lstlisting}
for _con in m.component_objects(pyo.Constraint):
    _keys_at_t0 = [k for k in _con.keys()
                   if (k[-1] if isinstance(k, tuple) else k) == t0]
    for k in _keys_at_t0:
        del _con[k]
\end{lstlisting}

This invariant is enforced by dedicated tests that build a finite or infinite block and verify that zero constraint entries exist at the relevant endpoint(s).

\subsection{Options Dataclass Design}

All configuration lives in a single \texttt{Options} dataclass with defaults for every field. This means:
\begin{itemize}[leftmargin=2em]
  \item Adding a new option never breaks existing code.
  \item The \texttt{for\_model\_module()} classmethod provides a natural entry point that applies model-level defaults, then caller overrides.
  \item \texttt{copy(**overrides)} supports immutable-style updates for creating variants.
\end{itemize}

Options grew substantially over the semester as new capabilities were added. Key additions include \texttt{objective} (replacing the old boolean \texttt{custom\_objective}), \texttt{lyap\_constraint\_type} (hard/soft/none), \texttt{terminal\_constraint\_type} (hard/soft/none), and \texttt{revive\_run} for automatic retry on near-convergent IPOPT failures.

\subsection{Lyapunov Stability Constraint}

The Lyapunov decrease condition is:
\[
  V_k - V_{k-1} \leq -\delta \cdot L_{k-1}
\]
where $V_k = \phi_\text{track}^\text{finite} + \beta \cdot \phi_\text{track}^\text{infinite}[\tau=1]$ is the combined quadratic tracking measure and $L_{k-1}$ is the first finite-element stage cost from the previous iteration.

Three constraint forms are supported via \texttt{lyap\_constraint\_type}:
\begin{itemize}[leftmargin=2em]
  \item \textbf{hard}: strict inequality constraint. Can be locally infeasible after a shift when the horizon is short.
  \item \textbf{soft}: the RHS gains a \texttt{lyap\_slack} variable (NonNegativeReals) that enters the objective with weight \texttt{lyap\_soft\_weight}. Recommended for short horizons.
  \item \textbf{none}: $\phi_\text{track}$ and $V_\text{prev}$ infrastructure is built but no constraint is added, allowing monitoring without enforcement.
\end{itemize}

\subsection{Solver Hardening}

Two reliability improvements were made after observing random failures at $\sim$90\% completion of long runs:

\textbf{Warm solver settings:} \texttt{bound\_relax\_factor=1e-8} (instead of 0) allows IPOPT to return solutions that are $10^{-8}$ outside a bound, which prevents the ``restoration phase converged to small primal infeasibility'' failure pattern caused by floating-point residuals introduced by \texttt{shift\_values\_by\_time}.

\textbf{Bound clipping:} After every successful solve, \texttt{\_clip\_to\_bounds(model)} projects all variable values back to $[\ell_b, u_b]$. This prevents downstream \texttt{set\_value()} failures on NonNegativeReals variables (e.g., setting a slack to $-10^{-9}$).

\textbf{Revive solver:} The \texttt{revive\_run} option (default 0) enables automatic retry with looser tolerances (\texttt{tol=1e-6}) when a warm solve fails with a retryable termination condition (not infeasible or maxIterations).

%% ============================================================
\section{The User Model Interface}
%% ============================================================

A user connects their process to the framework by writing a model file with two required functions:

\begin{lstlisting}
def variables_initialize(m):
    m.MV_index = pyo.Set(initialize=["Fa0"])
    m.CV_index = pyo.Set(initialize=["Ca", "Cb"])
    m.Ca = pyo.Var(m.time, initialize=0.5, domain=pyo.NonNegativeReals)
    m.Cb = pyo.Var(m.time, initialize=0.5, domain=pyo.NonNegativeReals)
    m.Fa0 = pyo.Var(m.time, initialize=12, bounds=(10, 20))
    m.dCadt = DerivativeVar(m.Ca, wrt=m.time)
    m.dCbdt = DerivativeVar(m.Cb, wrt=m.time)
    m.setpoints = pyo.Param(m.CV_index,
                            initialize={"Ca": 0.5, "Cb": 0.5})
    m.initial_values = pyo.Param(m.CV_index,
                                 initialize={"Ca": 0.0, "Cb": 0.0})
    return m

def equations_write(m):
    V, k, caf = 10, 1.2, 1
    m.Ca_balance = pyo.Constraint(m.time, rule=lambda m, t:
        m.dCadt[t] == m.Fa0[t]/V*(caf - m.Ca[t]) - k*m.Ca[t])
    m.Cb_balance = pyo.Constraint(m.time, rule=lambda m, t:
        m.dCbdt[t] == m.Fa0[t]/V*(0 - m.Cb[t]) + k*m.Ca[t])
    return m

def custom_objective(m, options):
    return lambda m, t: -m.Fa0[t]*(2*m.Cb[t] - 0.5)
\end{lstlisting}

The framework injects \texttt{m.time} before calling \texttt{variables\_initialize}. The model file must not declare \texttt{m.time}. After \texttt{equations\_write} runs, the framework deletes all constraint entries at non-collocation endpoints.

A built-in validator (\texttt{check-model path/to/model.py}) runs three levels of checks: static AST (syntax, required functions, no \texttt{m.time} declaration), import (module loads cleanly), and dynamic Pyomo (required components present after execution).

%% ============================================================
\section{Examples}
%% ============================================================

\subsection{Economic CSTR with Lyapunov Constraint}

The canonical test case is an isothermal CSTR with consecutive reaction A$\to$B:
\[
  \dot{C}_A = \frac{F_{A0}}{V}(C_{Af} - C_A) - kC_A, \qquad
  \dot{C}_B = -\frac{F_{A0}}{V}C_B + kC_A
\]
with manipulated variable $F_{A0} \in [10, 20]$ kmol/h and economic objective to maximize B production: $L(x,u) = -F_{A0}(2C_B - 0.5)$.

The controller was run for 100 MPC iterations with the following configuration:
\begin{center}
\begin{tabular}{ll}
\toprule
\textbf{Parameter} & \textbf{Value} \\
\midrule
Sampling time $T_s$ & 1 h \\
Finite elements (finite block) & 2, $n_\text{cp}=1$ (RADAU) \\
Finite elements (infinite block) & 5, $n_\text{cp}=1$ (LEGENDRE) \\
Auto-selected $\gamma$ & 0.100 h$^{-1}$ \\
Lyapunov constraint & hard, $\delta = 0.01$ \\
Average IPOPT solve time & 27 ms \\
\bottomrule
\end{tabular}
\end{center}

\begin{figure}[h]
  \centering
  \begin{subfigure}[b]{0.48\textwidth}
    \includegraphics[width=\textwidth]{figures/CV_1_C_a___mol_L_.png}
    \caption{$C_A$ trajectory (controlled variable)}
  \end{subfigure}
  \hfill
  \begin{subfigure}[b]{0.48\textwidth}
    \includegraphics[width=\textwidth]{figures/CV_2_C_b___mol_L_.png}
    \caption{$C_B$ trajectory (controlled variable)}
  \end{subfigure}
  \vspace{0.5em}
  \begin{subfigure}[b]{0.48\textwidth}
    \includegraphics[width=\textwidth]{figures/MV_1_F__A0____kmol_h_.png}
    \caption{$F_{A0}$ trajectory (manipulated variable)}
  \end{subfigure}
  \caption{CSTR eNMPC with infinite-horizon approximation and Lyapunov constraint. The system starts at $(C_A, C_B) = (0, 0)$ and converges to the economically optimal steady state. The Lyapunov constraint enforces monotone decrease of $V_k$ at each step.}
  \label{fig:cstr}
\end{figure}

The results show the system converging from the zero initial condition to the economic optimum. The infinite-horizon terminal cost provides foresight over the full future, while the Lyapunov constraint enforces a formal decrease certificate at each sample time.

\subsection{Ternary Distillation Column}

The ternary distillation example separates a three-component mixture A/B/C in two columns in series. It has 246 differential state variables (tray-by-tray liquid compositions and holdups), 8 manipulated variables (vapor boilup and liquid reflux for each column, plus distillate and bottoms flows), and 3 algebraic controlled variables (distillate purities $x_{D1A}$, $x_{D2B}$, and bottoms composition $x_C$). The economic objective minimizes feed and energy cost while maximizing product revenue.

This example demonstrates the framework's ability to handle large-scale DAE systems (246 differential states, hundreds of algebraic variables) while maintaining the full infinite-horizon structure. Soft Lyapunov and terminal constraints are used because the short finite horizon makes hard constraints locally infeasible after each shift.

%% ============================================================
\section{Testing}
%% ============================================================

The test suite contains 121 tests across 7 files, run via GitHub Actions CI on every push to \texttt{main}. The suite was designed so that each test documents \emph{why} an invariant exists, not just that it holds.

\begin{center}
\begin{tabular}{lll}
\toprule
\textbf{File} & \textbf{Tests} & \textbf{Coverage} \\
\midrule
\texttt{test\_options.py} & 28 & Dataclass defaults, classmethods, copy semantics \\
\texttt{test\_check\_model.py} & 12 & Static/dynamic model validator, pass/fail/warn cases \\
\texttt{test\_collocation\_tools.py} & 10 & $\gamma$ formula, Lagrange coefficients (partition of unity) \\
\texttt{test\_solver\_tools.py} & 9 & \texttt{\_clip\_to\_bounds}: below lb, above ub, nested blocks \\
\texttt{test\_model\_structure.py} & 13 & t=0 deletion (finite block), $\tau=0/1$ deletion (infinite block) \\
\texttt{test\_controllers.py} & 29 & Block structure, lyap types, terminal types, plant MV fixing \\
\texttt{test\_mpc\_loop.py} & 20 & End-to-end convergence, CSV output, six controller variants \\
\bottomrule
\end{tabular}
\end{center}

The structural tests (model structure, endpoint deletion) are particularly valuable: they run in $\sim$2 seconds without IPOPT and would immediately catch any regression in the $t=0$ deletion logic that was the root cause of the most damaging bug in the project's history.

%% ============================================================
\section{Agentic Engineering: Analysis and Reflection}
%% ============================================================

\subsection{What Agentic Engineering Meant in Practice}

This project was developed almost entirely through agentic engineering with Claude. By ``agentic'' I mean that the AI was given not just isolated editing tasks, but persistent context about the project state, freedom to explore codebases, write multiple files at once, run tests, read error messages, and iterate---functioning more like a pair programmer than a code-completion tool.

The two most important artifacts that made this work were:

\textbf{LLMREADME.md} --- a living technical reference document describing the entire codebase architecture, key invariants, and conventions. Unlike a normal README (written for human newcomers), this document was written in a way that a future AI agent could read and immediately understand \emph{why} things are the way they are, not just what they are. It was updated in every major session.

\textbf{LLMJOURNAL.md} --- a chronological log of every session: what problem was being worked on, what the root cause turned out to be, what was changed, and what was left unresolved. This served as the ``conversation history'' across sessions that a stateless AI would otherwise lose.

\subsection{How the Experience Evolved}

\textbf{Early semester: isolated edits, frequent context loss.} The first few sessions involved asking Claude to make specific changes to specific functions. Each session started from scratch, requiring substantial re-explanation of the project. The AI would often suggest changes that conflicted with constraints it hadn't been told about (e.g., the \texttt{clean\_model='delete'} requirement, which ruled out many standard Pyomo patterns).

The core insight from this phase was: \emph{context is the bottleneck}. The AI is capable, but only within the scope of what it knows. A session that starts with ``fix the warm-start bug'' and no other context produces a different (and usually worse) answer than one that starts with ``here is the full architecture, here is the IPOPT output, here is what we've already tried.''

\textbf{Mid semester: structured context and longer sessions.} Once LLMREADME and LLMJOURNAL were established, sessions became dramatically more productive. A single session in April addressed four interconnected problems (plant infeasibility, load\_data skipping fixed variables, Lyapunov constraint violations, equations\_write at $t=0$) and produced a correct, tested fix for each. The AI could hold enough project context to reason about interactions between components rather than treating each file in isolation.

The key workflow that emerged was:
\begin{enumerate}[leftmargin=2em]
  \item Open a session by pointing the AI to LLMREADME and LLMJOURNAL.
  \item Describe the observed failure and any diagnostic output already collected.
  \item Let the AI read the relevant source files and propose a fix with explicit reasoning.
  \item Run the fix and report results. If it fails, share the new error and iterate.
  \item At the end of the session, ask the AI to update LLMJOURNAL with a summary.
\end{enumerate}

\textbf{Late semester: code generation at scale.} The final sessions (test suite, CI, report) involved generating large amounts of new code rather than debugging existing code. Here agentic engineering excelled: the AI could read the entire LLMREADME, understand what invariants existed, and write 121 tests that checked exactly those invariants with accurate docstrings explaining why each test existed.

This would have taken far longer manually---not because the tests are hard to write, but because the mental overhead of reading the architecture, mapping invariants to test cases, and writing clear docstrings for 121 tests is substantial. The AI did it in one session.

\subsection{What Worked Well}

\begin{itemize}[leftmargin=2em]
  \item \textbf{Debugging with full context.} Given IPOPT output, a model pprint, and a description of what was expected, the AI consistently identified root causes correctly. The $t=0$ deletion bug (root cause 4 of the plant infeasibility) was found by the AI after a few iterations of providing diagnostic outputs.
  \item \textbf{Refactoring with invariants.} When the \texttt{custom\_objective} bool was replaced with the \texttt{objective}/\texttt{tracking\_setpoint} string pair, the AI correctly propagated the change across 9 example files, 3 modules, and the documentation in a single session.
  \item \textbf{Documentation generation.} The LLMREADME is more thorough and accurate than any documentation I would have written manually, because the AI simultaneously read the code and wrote the description---no translation step.
\end{itemize}

\subsection{What Required Care}

\begin{itemize}[leftmargin=2em]
  \item \textbf{The AI cannot test live.} When a proposed fix required running IPOPT (a 30-second simulation), the iteration cycle was manual: apply the fix, run the example, share the result. This is inherent to any AI-assisted workflow with external solvers.
  \item \textbf{Hallucination of non-existent APIs.} Early in the project, the AI occasionally suggested Pyomo functions that did not exist in the forked version, or that existed in the standard version but were incompatible with \texttt{clean\_model='delete'}. The fix was to ensure LLMREADME documented which Pyomo patterns were allowed, and to always verify suggestions against the actual source.
  \item \textbf{Scope creep.} Without explicit scope constraints, the AI sometimes refactored more than requested (e.g., cleaning up style while fixing a bug). This was addressed by adding ``don't refactor beyond what the task requires'' as an explicit instruction in the project's CLAUDE.md.
\end{itemize}

\subsection{Overall Assessment}

Agentic engineering made this project possible at its current scale. A 246-state ternary distillation model with Lyapunov constraints, an automatic time-compression parameter, three constraint-type options, a 121-test suite, and CI---built by one person in one semester---would not have been feasible without an AI pair programmer that could hold the full codebase in context.

The most important skill was learning to treat the AI as a collaborator that needs to be briefed, not a tool that just runs commands. Investing time in LLMREADME and LLMJOURNAL returned compounding dividends: each session was faster and produced higher-quality output than the last, because the AI could immediately understand the constraints and history rather than rediscovering them.

%% ============================================================
\section{Conclusion}
%% ============================================================

\textbf{infeNMPC} is a complete, tested, and documented Python framework for infinite-horizon economic NMPC. It solves the problem of integrating the full cost-to-go into a finite optimization problem via the tanh time transformation, with an optional Lyapunov stability certificate. The user interface---a model file with two required functions---is simple enough for a new user to adopt within a day while being expressive enough to handle 246-state industrial-scale DAE systems.

The project was developed entirely through agentic engineering. The central lesson is that structured context artifacts (LLMREADME, LLMJOURNAL) are as important as the AI model itself. A capable AI with poor context produces mediocre results; the same AI with rich, accurate context produces production-quality code with minimal iteration.

\end{document}
